\documentclass[11pt,a4paper]{article}
\usepackage[margin=25mm]{geometry}
\usepackage{amsmath,amssymb,amsthm,booktabs,hyperref}
\newtheorem{theorem}{Theorem}
\newtheorem{lemma}[theorem]{Lemma}
\newcommand{\Sz}{\operatorname{Sz}}
\newcommand{\gen}[1]{\langle#1\rangle}
\newcommand{\F}{\mathbb F}
\title{A chiral $4$-polytope with automorphism group $\Sz(32768)$}
\author{Human Compression project}
\date{}
\begin{document}\maketitle
\begin{abstract}
We give explicit matrices for a chiral polytope of type $\{31,31,25\}$
whose automorphism group is $\Sz(2^{15})$. Its facet and vertex-figure
groups are two conjugate subfield groups $\Sz(2^5)$, meeting in a split
torus of order $31$. The construction contradicts Theorem~1.1 in the
specified preprint edition of Spiga's paper on Suzuki polytopes.
All finite-field calculations are supplied in a short executable
certificate. The proof of generation and of the intersection property
uses the classical structure of Suzuki groups, not an enumeration of
the ambient group.
\end{abstract}

Spiga's preprint~\cite{Spiga}, arXiv:2010.04080v1 of 8 October 2020,
asserts that every chiral polytope whose automorphism group lies between
a Suzuki simple group and its automorphism group has rank three.
The question is attributed there to Hubard and Leemans~\cite{HL}.
The theorem below concerns that precise source edition. It does not
assert that every later printing has the same text, or make a separate
priority claim about the subsequent literature.

\begin{theorem}\label{main}
There exists a chiral abstract polytope of rank four and type
$\{31,31,25\}$ with automorphism group $\Sz(32768)$.
\end{theorem}

\section{Coordinates and retained group structure}
Let $k=\F_{2^{15}}$ and $\theta=2^8$. We use the standard Suzuki
matrices
\[
 U(x,y)=\begin{pmatrix}
 1&0&0&0\\x&1&0&0\\x^{\theta+1}+y&x^\theta&1&0\\
 x^{\theta+2}+xy+y^\theta&y&x&1
 \end{pmatrix},\qquad
 H(t)=\operatorname{diag}(t^{\theta+1},t,t^{-1},t^{-\theta-1}),
\]
where $x,y\in k$ and $t\in k^*$. Let $J$ be the matrix with ones on
the antidiagonal and zeroes elsewhere. These matrices generate
$G=\Sz(k)$; see, for example, the coordinate description in
Bray and B\"a\"arnhielm~\cite{Recognition}. Write $K=\F_{32}$ and
$F=\Sz(K)$ for the standard subfield group. The restriction of the
map $x\mapsto x^\theta$ to $K$ is $x\mapsto x^8$, so the same
formulas define $F$.

We use the following classical facts from Suzuki's subgroup
classification~\cite{Suzuki}. The group $\Sz(q)$ is simple for
$q\geq8$, has order $q^2(q^2+1)(q-1)$, and acts doubly transitively
on its ovoid. A nonidentity element of a split torus fixes exactly the
two points defining that torus. The maximal subgroups are point
stabilizers, the three types of odd-torus normalizers, and maximal
Suzuki subfield groups. Their relevant orders are
\[
 q^2(q-1),\qquad 2(q-1),\qquad
 4(q+\sqrt{2q}+1),\qquad4(q-\sqrt{2q}+1).
\]
A subfield group $\Sz(2^d)$ is maximal when the odd quotient of field
exponents is prime. The centralizer of a nonidentity odd-order element
is cyclic of odd order. Finally,
$\operatorname{Out}(\Sz(2^N))\cong C_N$ for odd $N$.
These are borrowed structural inputs. None is inferred from the
finite computation below.

The standard Bruhat decomposition, also retained, is
\[
 \Sz(K)=U(K)H(K)\ \sqcup\ U(K)H(K)J U(K).
 \tag{1}
\]
The first cell is lower triangular. In the second cell the entry in
row one, column four is nonzero. Both assertions follow immediately
by multiplying the displayed matrices.

\section{The explicit rotations}
Take $z$ to be the residue class of $X$ in
\[
 k=\F_2[X]/(X^{15}+X+1),\qquad \eta=z^{1057}.
 \tag{2}
\]
The polynomial is irreducible, $z$ has order $32767$, and $\eta$ has
order $31$. Thus $K=\F_2(\eta)$. Set
\[
 \lambda=z^{11}+z^7+z^6+z^4+z^3+z+1.
 \tag{3}
\]
For reference, its minimal polynomial over $K$ is
\[
 T^3+\eta^{12}T^2+\eta^{11}T+\eta.
 \tag{4}
\]
In particular $\lambda\notin K$.
Define the involution
\[
 i(x,y)=U(x,y)^{-1}J U(0,1)J U(x,y)
\]
and put
\[
 \begin{aligned}
 u&=i(\eta^2,\eta^3),&v_0&=i(\eta^5,\eta^{30}),\\
 h&=H(\lambda),&v&=h v_0h^{-1},&b&=H(\eta),\\
 a&=u b^{-1},&&c=b^{-1}v.
 \end{aligned}
 \tag{5}
\]
These formulas, together with (2)--(3), specify every entry without a
choice of a primitive element or of an unspecified root.

\begin{lemma}\label{arithmetic}
The elements in (5) satisfy
\[
 |a|=31,\quad |b|=31,\quad |c|=25,\qquad
 |ab|=|bc|=|abc|=2.
 \tag{6}
\]
Neither $u$ nor $v_0$ is upper triangular, lower triangular, or
antidiagonal.
\end{lemma}
\begin{proof}
The assertions are exact calculations in the field (2). We give the
calculation in executable form as \texttt{check\_arithmetic.py}; it
uses only polynomial arithmetic over $\F_2$ and multiplication of
$4\times4$ matrices. An integer $m=\sum m_j2^j$ denotes
$\sum m_jz^j$. In this notation $\eta=624$ and $\lambda=2267$,
and the two involutions occurring before conjugation are
\[
 u=\begin{pmatrix}
 4671&4175&5401&1\\5156&1562&1830&5401\\
 845&5157&1562&4175\\1131&845&5156&4671
 \end{pmatrix},\quad
 v_0=\begin{pmatrix}
 5401&5401&4867&1\\4467&1563&1562&4867\\
 1131&5717&1563&5401\\624&1131&4467&5401
 \end{pmatrix}.
 \tag{7}
\]
Reduction modulo $X^{15}+X+1$ gives $a^{31}=b^{31}=c^{25}=I$,
with $a,b,c^5\ne I$. It also gives
$u^2=v^2=(abc)^2=I$, with all three matrices nonidentity.
The identities $ab=u$ and $bc=v$ follow from the definitions.
The off-diagonal and diagonal entries in (7) establish the last
assertion directly.

For reproducibility, the certificate first checks irreducibility by
$z^{2^{15}}=z$ and the two polynomial gcd tests for the prime divisors
$3,5$ of $15$. It verifies primitivity by the prime divisors
$7,31,151$ of $32767$, reconstructs every matrix from (2)--(5), and
checks all the equalities just stated. It also verifies the Suzuki
and symplectic matrix equations and rejects a deliberately altered
entry of $v$. The earlier search output is not read by this checker.
\end{proof}

\section{The two face groups and their intersection}
\begin{lemma}\label{faces}
We have $\gen{a,b}=F$ and $\gen{b,c}=hFh^{-1}$.
\end{lemma}
\begin{proof}
The first group equals $\gen{u,b}\leq F$. At $q=32$, the two
nonsplit torus normalizers have orders $100$ and $164$, and the
proper Suzuki subfield group is $\Sz(2)$, of order $20$. None has
order divisible by $31$. A proper subgroup containing $b$ must
therefore lie in a point stabilizer or in the normalizer of a split
torus. The only ovoid points fixed by $b$ are the standard points
$0,\infty$. Their stabilizers are upper and lower triangular,
respectively. The normalizer of $\gen b=H(K)$ consists of the
diagonal matrices and their antidiagonal coset. Lemma~\ref{arithmetic}
excludes every possibility for a proper subgroup containing $u,b$.
This proves the first assertion.

The element $h$ centralizes $b$. Consequently
$h^{-1}\gen{b,c}h=\gen{b,v_0}$. The same argument, with $v_0$ in
place of $u$, proves the second assertion.
\end{proof}

\begin{lemma}\label{intersection}
For the diagonal conjugator in (5),
\[
 F\cap hFh^{-1}=H(K)=\gen b.
 \tag{8}
\]
\end{lemma}
\begin{proof}
The diagonal subgroup $H(K)$ is contained in the intersection.
Suppose $g\in F\cap hFh^{-1}$. Then both $g$ and $h^{-1}gh$
have entries in $K$. If $g$ belongs to the second Bruhat cell in
(1), its $(1,4)$ entry is in $K^*$, whereas conjugation by $h$
multiplies this entry by $\lambda^{-2(\theta+1)}$. Hence
$\lambda^{514}\in K$. Since $\gcd(514,32767)=1$, this implies
$\lambda\in K$, a contradiction.

Thus $g=U(x,y)H(t)$ with $x,y\in K$ and $t\in K^*$. Direct
conjugation gives
\[
 h^{-1}U(x,y)h
   =U(x\lambda^\theta,y\lambda^{\theta+2}).
 \tag{9}
\]
The root parameters in (1) are unique, so the two parameters on the
right belong to $K$. If $x\ne0$, then $\lambda^{256}\in K$;
if $x=0$ and $y\ne0$, then $\lambda^{258}\in K$.
Both are impossible because
$\gcd(256,32767)=\gcd(258,32767)=1$. Therefore $x=y=0$ and
$g\in H(K)$, as required.
\end{proof}

\begin{lemma}\label{generation}
The three rotations generate $G=\Sz(32768)$.
\end{lemma}
\begin{proof}
They generate both face groups from Lemma~\ref{faces}. The first,
$F=\Sz(32)$, is a maximal subgroup of $G$ because the quotient of
field exponents is $15/5=3$. Equation (8) implies that the second
face group is distinct from $F$. Hence the generated subgroup
properly contains $F$ and must equal $G$.
\end{proof}

\section{The polytope and its chirality}
We recall the rank-four rotation criterion of Schulte and
Weiss~\cite{SW}, in the explicit form recorded in
Leemans, Moerenhout and O'Reilly-Regueiro~\cite[\S2, (C3$'$)]{LMO}.
Given generators $a,b,c$ of finite orders at least two, with
$ab,bc,abc$ involutions, the required intersection condition is
\[
 \gen a\cap\gen b=1,\qquad
 \gen b\cap\gen c=1,\qquad
 \gen{a,b}\cap\gen{b,c}=\gen b.
 \tag{10}
\]
These conditions give a chiral or directly regular abstract polytope.
It is chiral precisely when there is no involutory automorphism
$\alpha$ of the generated group satisfying
\[
 \alpha(a)=a^{-1},\qquad \alpha(b)=a^2b,
 \qquad\alpha(c)=c.
 \tag{11}
\]
This is a borrowed reconstruction theorem; the group and marking
constructed in (5) are the new input.

\begin{proof}[Proof of Theorem~\ref{main}]
The rotation relations and generation follow from
Lemmas~\ref{arithmetic} and \ref{generation}. The groups $\gen a$
and $\gen b$ have prime order $31$. If they had nontrivial
intersection, they would coincide, and their product $ab$ would
have odd order, contrary to (6). Thus the first equality in (10)
holds. The second follows from $\gcd(31,25)=1$, and the third is
Lemma~\ref{intersection}.

It remains to exclude (11). The outer automorphism group of $G$ has
odd order $15$, so every involutory automorphism is inner. Write
$\alpha(g)=s^{-1}gs$. Since $G$ is centreless, $\alpha^2=1$
implies $s^2=1$. The automorphism cannot be the identity, because
$|a|=31$, so $s$ would be an involution. But $\alpha(c)=c$ gives
$s\in C_G(c)$, and this centralizer has odd order because $|c|=25$.
This is impossible. The polytope is therefore chiral, with full
automorphism group $G$ and type $\{31,31,25\}$.
\end{proof}

The computation also checks the two small dihedral groups
$\gen{a,bc}$ and $\gen{ab,c}$, of orders $62$ and $50$, and their
intersections with the face groups. Their mutual intersection has
order two and is generated by $abc$. These finite checks corroborate
(10); they do not replace the proof of (8) or the maximal-subgroup
argument. The ambient group, of order
$37\,777\,778\,976\,635\,853\,209\,600$, is never enumerated.

\section{Comparison with the source coordinates}
For a direct comparison with~\cite{Spiga}, let $P$ be the permutation
matrix for $(1\ 2)(3\ 4)$. It preserves $J$ and changes our torus to
\[
 PH(t)P^{-1}
  =\operatorname{diag}(t,t^{\theta+1},t^{-\theta-1},t^{-1}).
\]
Our Suzuki graph map takes the minors with row and column pair order
$(12,13,24,34)$, then raises entries to $2^7$. Conjugating by $P$
changes that order to $(24,12,34,13)$, exactly the order in Wilson's
product as displayed on page one of~\cite{Spiga}. This verifies the
change of model rather than assuming identical coordinate conventions.

The certificate constructs $g\in G$ with $g^{-1}(abc)g=J$.
In the source basis, the inverse of the first normalized rotation is
\[
 x=P(g^{-1}ag)^{-1}P^{-1}
  =\begin{pmatrix}
 29124&17285&29707&20757\\
 17285&3002&31396&7939\\
 29707&31396&8273&19487\\
 20757&7939&19487&23604
 \end{pmatrix}.
 \tag{12}
\]
The integer encoding is the one used in (7). This matrix is symmetric,
belongs to the Suzuki group in the source model, and satisfies
$xJx=J$. Nevertheless its entries give
\[
 x_{14}x_{24}+x_{12}x_{44}=28331\ne0.
 \tag{13}
\]
Thus it violates equation (13) of~\cite{Spiga}, whose right-hand side
would be zero.

The reason is visible before that equation is imposed. The source
states Wilson's condition $gx\mathbin\bullet gy=g(x\mathbin\bullet y)$
only for pairs with symplectic pairing zero. On page three it also
imposes this condition on $(e_1,e_4)$ and $(e_2,e_3)$, both of which
have pairing one. For the concrete matrix (12),
\[
 (xe_1)\mathbin\bullet(xe_4)
   =(3245,26415,3338,22277)^{\mathsf T}\ne0,
 \qquad e_1\mathbin\bullet e_4=0.
\]
The admissible orthogonal-pair identities do hold. The failure is
therefore an application outside the stated domain of the product
identity, not a discrepancy between the two matrix models.

There is also a count that detects the same defective classification.
For $\mathcal X=\{x\in G:xJx=J\}$, the map $x\mapsto xJ$ is a
bijection from $\mathcal X$ to the involutions together with the
identity. Each of the $q^2+1$ ovoid points is fixed by $q-1$
involutions, and each involution has a unique fixed point. Hence
$|\mathcal X|=(q^2+1)(q-1)+1$, whereas the classification in the source
would give $|\mathcal X|=q$. This observation diagnoses the source
argument; Theorem~\ref{main} additionally verifies the polytope
intersection and chirality conditions.

The contribution of this note is the explicit marking, the diagonal
subfield intersection argument as applied to it, and the resulting
counterexample to the stated rank bound. The matrix model, Suzuki
subgroup classification and polytope reconstruction theorem are
retained inputs with the attributions above. No whole-source
acceptance, independent review, Lean verification, or compression
measurement is claimed by this author packet itself.

\begin{thebibliography}{6}
\bibitem{Spiga} P. Spiga, \emph{On the rank of Suzuki polytopes:
an answer to Hubard and Leemans}, arXiv:2010.04080v1,
8 October 2020, Theorem 1.1.
\bibitem{HL} I. Hubard and D. Leemans, \emph{Chiral polytopes and
Suzuki simple groups}, Rigidity and Symmetry, Fields Institute
Communications 70 (2014), 155--175.
\bibitem{Suzuki} M. Suzuki, \emph{On a class of doubly transitive
groups}, Annals of Mathematics 75 (1962), 105--145.
\bibitem{Recognition} J. N. Bray and H. B\"a\"arnhielm,
\emph{A new method for recognising Suzuki groups},
arXiv:1706.05697v1 (2017), \S2.
\bibitem{SW} E. Schulte and A. Ivi\'c Weiss, \emph{Chiral polytopes},
Applied Geometry and Discrete Mathematics, DIMACS Series 4 (1991),
493--516.
\bibitem{LMO} D. Leemans, J. Moerenhout and E. O'Reilly-Regueiro,
\emph{Projective linear groups as automorphism groups of chiral
polytopes}, arXiv:1606.08017v1, 26 June 2016, \S2, printed pages 3--4.
\end{thebibliography}
\end{document}
